\section{限制與未來研究}

本研究首先受限於硬體條件。雖然 ModernBERT 原生支援 8192 tokens，但在 \texttt{query:positive:negatives=1:1:15} 的對比式學習配置下，單張 RTX 4080 SUPER 16GB 無法穩定訓練 8192-token 版本，因此本研究最終採用 4096 tokens。這意味著目前結果仍未完全驗證「真正的 8192-token 對比式法律檢索」是否能帶來更大提升。未來若能使用更大顯存的 GPU、分散式訓練、gradient checkpointing 或更高效的 chunking/aggregation 設計，將能更完整檢驗長文本假設。

第二，法律領域 continued pretraining 目前僅訓練約 60 小時，仍只覆蓋不到半個 epoch，尚未完整看過全部 6,735,525 篇美國判決書，因此其潛力很可能尚未完全釋放。未來可延長 DAP 步數，或比較美國、加拿大與混合法域語料對 COLIEE 表現的影響。

第三，雖然本文已經把 LightGBM learning-to-rank 與 cut-off post-processing 納入正式提交流程，並在官方榜單上帶來可觀但有限的提升，但目前的 reranking 仍主要依賴 bi-encoder 產生的特徵與 query-wise tree ranker。未來若能進一步結合 cross-encoder、listwise reranker，或更細緻的 chunk interaction 機制，理論上仍有持續提升空間。

第四，本文雖然已驗證 year-based scope filtering 有助於法律檢索，但目前的年份抽取仍主要依賴規則與 metadata/文本模式，未來可以把年份、程序階段與引用關係抽取得更完整，讓候選範圍約束不只限於 ``是否早於 query''，而能進一步納入法域、法院層級與程序類型等法律知識。

最後，從工程角度來看，本研究觀察到 fp16 fine-tuning 比 bf16 更容易收斂，但其真正原因尚未完全釐清。後續可系統性比較不同 precision、learning rate、temperature 與 hard-negative 更新策略，釐清數值穩定性與法律檢索表現之間的關係。
